\section{Methods}
\label{sec:methods}

\subsection{Datasets}
\label{sec:datasets}

We assembled five publicly available resting-state EEG datasets comprising 2,097 participants and nine eyes-closed (EC) recording conditions (Table~\ref{tab:datasets}).

\begin{table}[htbp]
\centering
\caption{EEG datasets included in the analysis. N = number of subjects contributing at least one valid EC resting-state recording. HBN releases R1--R6 (R5 excluded due to artefact contamination) are treated as independent replications within the paediatric cohort.}
\label{tab:datasets}
\begin{tabular}{llrrll}
\toprule
\textbf{Dataset} & \textbf{N} & \textbf{Channels} & \textbf{Age (years)} & \textbf{Country} & \textbf{Conditions} \\
\midrule
EEGMMIDB & 109 & 64 & adult & USA & 14 runs pooled \\
LEMON & 203 & 59 & 20--77 & Germany & EC, EO \\
Dortmund & 608 & 64 & 18--70 & Germany & EC/EO $\times$ pre/post $\times$ ses-1/ses-2 \\
CHBMP & 250 & 62--120 & adult & Cuba & EC \\
HBN R1--R6 & 927 & 128 & 5--21 & USA & Mixed resting \\
\midrule
\textbf{Total} & \textbf{2,097} & & & & \textbf{9 EC conditions} \\
\bottomrule
\end{tabular}
\end{table}

The \textbf{EEG Motor Movement/Imagery Database (EEGMMIDB)} comprises 109 healthy adults recorded with a 64-channel system at 160 Hz \parencite{goldberger2000physiobank,schalk2004bci2000}. Subjects performed fourteen one-minute runs covering rest and motor imagery tasks; all runs were pooled per subject.

The \textbf{LEMON} (Leipzig Study for Mind-Body-Emotion Interactions) dataset comprises 203 healthy adults (ages 20--77) recorded with 59 channels at 2500 Hz \parencite{babayan2019lemon}. Paired EC and EO resting-state blocks ($\approx$15 minutes each) and a neuropsychological battery including the Leistungsprüfsystem (LPS), Test of Attentional Performance (TAP), Regensburger Wortflüssigkeits-Test (RWT), and Trail Making Test (TMT) make LEMON the primary cognitive-correlates dataset.

The \textbf{Dortmund EEG dataset} is a longitudinal sample of 608 healthy German adults (ages 18--70) recorded at 64 channels, with a subset of $N = 208$ returning approximately five years later \parencite{vossen2015dortmund}. The longitudinal subsample provides five-year test--retest reliability estimates.

The \textbf{Cuban Human Brain Mapping Project (CHBMP)} provides EC resting-state recordings from 250 healthy Cuban adults recorded at 62--120 channels.

The \textbf{Healthy Brain Network (HBN)} dataset \parencite{alexander2017hbn} is a paediatric community sample of 927 participants (ages 5--21) from the New York City area, recorded at 128 channels. Five independent data releases (R1, R2, R3, R4, R6; R5 excluded due to artefact contamination) were analysed to test cross-sample reproducibility. HBN includes the Child Behavior Checklist (CBCL) for psychopathology measures.

\subsection{EEG Preprocessing and Spectral Peak Extraction}
\label{sec:extraction}

All datasets were processed through a common v3 pipeline. Raw recordings were bandpass filtered from 1 to 59~Hz with a zero-phase fourth-order Butterworth filter. European recordings (LEMON and Dortmund) received an additional 50~Hz notch filter ($Q = 30$) to suppress mains-frequency contamination. All recordings were resampled to 250~Hz.

Spectral peaks were extracted using the FOOOF/specparam framework \parencite{donoghue2020fooof}, which models each power spectrum as the sum of an aperiodic $1/f$-like component and Gaussian peaks. Adaptive Welch windowing with a minimum segment length of 2500 samples (10~s at 250~Hz) yielded frequency resolution $\Delta f = 0.1$~Hz. The bandwidth floor was set to $2 \times \Delta f = 0.2$~Hz.

Four fitting ranges were applied sequentially:
\begin{enumerate}
    \item \textbf{Merged theta--alpha}: 4.70--12.30~Hz, single aperiodic component.
    \item \textbf{Beta-low}: 12.30--19.90~Hz.
    \item \textbf{Beta-high}: 19.90--32.19~Hz.
    \item \textbf{Gamma}: 32.19--52.09~Hz.
\end{enumerate}

Global specparam parameters: \texttt{peak\_threshold} $= 0.001$, \texttt{max\_n\_peaks} $= 15$, \texttt{min\_peak\_height} $= 0.0001$. Each peak was retained only if $R^2 \geq 0.70$ (median $R^2 = 0.931$ across all retained peaks). Each peak was characterised by centre frequency, power, bandwidth, $R^2$, and $\phisym$-octave assignment. After power filtering (Section~\ref{sec:power_filter}), the final dataset comprised 4,572,489 peaks from an initial 9,132,221.

\subsubsection{Three Methodological Corrections}
\label{sec:methods_corrections}

\paragraph{Correction 1: $\fzero$ alignment.}
Peaks are assigned to bands via the normalised coordinate:
\begin{equation}
u = \left(\frac{\log(f / \fzero)}{\log \phisym}\right) \bmod 1,
\label{eq:u_coord}
\end{equation}
where $\fzero = 7.60$~Hz. A prior 3\% mismatch between extraction ($\fzero = 7.83$~Hz) and enrichment coordinates ($\fzero = 7.60$~Hz) caused boundary-wrapping artifacts. Aligning both to $\fzero = 7.60$~Hz eliminated five sign reversals.

\paragraph{Correction 2: Hz-weighted null.}
Under a Hz-uniform detection process, the expected count in Voronoi bin $i$ with $u$-edges $[u_L, u_R]$ is:
\begin{equation}
\hat{E}_i = \frac{\phisym^{u_R} - \phisym^{u_L}}{\phisym - 1} \cdot N,
\label{eq:hz_weighted_null}
\end{equation}
where $N$ is total peaks in the band. Enrichment is:
\begin{equation}
\text{Enr}_i = \left(\frac{n_i}{\hat{E}_i} - 1\right) \times 100\%.
\label{eq:enrichment}
\end{equation}

\paragraph{Correction 3: Merged theta--alpha fit.}
Separate FOOOF fits for theta and alpha created a 15$\times$ step discontinuity in peak density at $\fzero$. Merging into a single fit (4.70--12.30~Hz) eliminated this artifact ($r = 0.913$ with separate alpha profile; $r = 0.200$ with separate theta, confirming the theta profile was dominated by the boundary artifact).

\subsubsection{Power Filtering}
\label{sec:power_filter}

Permissive extraction (\texttt{peak\_threshold} $= 0.001$) retained virtually all supra-background features. Analysis-time filtering then kept the top 50\% of peaks by power within each band per subject. This ``permissive extraction + analysis-time filtering'' strategy outperformed strict-threshold extraction on all validation metrics.

\subsubsection{Parameter Validation}
\label{sec:param_validation}

A 48-configuration parameter sweep on EEGMMIDB (GCP c2d-standard-32, $\sim$30 minutes total) factorially combined \texttt{peak\_threshold} $\in \{0.001, 0.01, 0.1, 2.0\}$, \texttt{min\_peak\_height} $\in \{0.00, 0.05, 0.10, 0.15, 0.20\}$, \texttt{max\_n\_peaks} $\in \{5, 8, 12, 15, 20\}$, and adaptive versus fixed Welch segmentation. Only \texttt{max\_n\_peaks} materially affected enrichment (inflection at 12, matching the 12 Voronoi positions). The production value of 15 was chosen for stable per-subject estimates.

\subsection{Log-Frequency Scaling Tests}
\label{sec:scaling_methods}

Seven tests evaluated whether the peak distribution is better described in log-frequency or linear-frequency space, and whether the inter-trough ratio is robust to the dependence structure of the data.

\textbf{Test 1: KDE density estimation.} Peak frequencies were pooled ($N = 4{,}572{,}489$) and Gaussian KDE applied in both linear-Hz and log-Hz spaces ($h = 0.02$). Local minima (density troughs) were identified as empirical band boundaries.

\textbf{Test 2: BIC model comparison.} Empirical boundary and centre frequencies were fitted by geometric series $f_k = \fzero \times r^k$ for seven fixed ratios ($\phisym$, $e-1$, $\sqrt{3}$, 2, $e$, $2^{1/3}$, $\sqrt{2}$), a free-ratio model (2 parameters), and linear-Hz spacing (2 parameters). BIC $= n \ln(\text{SSE}/n) + k \ln n$.

\textbf{Test 3: Polynomial BIC.} Peak frequencies were histogrammed in 100 equal-width bins in linear-Hz and log-Hz. Polynomial BIC (degrees 2--6) compared description parsimony.

\textbf{Test 4: Per-dataset replication.} Tests 1--3 were repeated independently on each of 9 datasets. One-sample $t$-tests on log-transformed ratios tested $H_0: \mu_{\log r} = \log r_\text{candidate}$.

\textbf{Test 5: Smoothing stability.} Trough locations were tracked across 30 smoothing bandwidths ($\sigma = 2$--30 bins on a 1000-bin log-frequency histogram). A trough was ``stable'' if present in $>50\%$ of smoothings.

\textbf{Test 6: Aperiodic null.} 200 Poisson surrogates were drawn from a heavily smoothed ($\sigma = 40$ bins) version of the real log-frequency histogram, preserving the $1/f$ envelope while destroying band structure. Trough depths were compared between real data and surrogates.

\textbf{Test 7: Subject-level bootstrap.} To account for the non-independence of peaks nested within subjects, 1,000 bootstrap resamples were drawn at the \emph{subject} level, stratified by dataset (preserving the 9-dataset composition at each iteration). For each resample, the full trough-detection pipeline (histogram $\to$ smoothing $\to$ peak-finding) was re-run from scratch on the pooled resampled peaks. Trough positions were matched to reference troughs algorithmically by nearest-neighbour assignment in log-frequency space (maximum distance: 0.3 log-Hz, corresponding to $\approx$35\% in Hz). The global inter-trough geometric mean ratio was computed for each bootstrap iteration. Confidence intervals were derived as the 2.5th and 97.5th percentiles. Two-sided bootstrap $p$-values for candidate constants were computed as $p = 2 \times \min(F(\theta), 1 - F(\theta))$, where $F(\theta)$ is the proportion of bootstrap replicates at or below the candidate value $\theta$.

\subsection{Boundary Sweep}
\label{sec:sweep_methods}

A $36 \times 36$ grid of $(\fzero, r)$ pairs spanning $\fzero \in [6.0, 9.5]$~Hz and $r \in [1.3, 2.3]$ was evaluated. For each coordinate system, band boundaries were placed at $\fzero \cdot r^n$. The normalised intra-band coordinate is:
\begin{equation}
u = \frac{\log(f / f_\text{lower})}{\log r},
\end{equation}
which is invariant to the choice of logarithm base. Within each band, peaks were binned into 12 equal-width $u$-bins with Hz-weighted expected counts (Equation~\ref{eq:hz_weighted_null}).

Three metrics were evaluated: \textbf{profile simplicity} (mean $R^2$ of best linear/quadratic fit across bands), \textbf{cross-dataset consistency} (mean pairwise Pearson $r$ of enrichment vectors across 9 datasets), and \textbf{band independence} (mean $|r|$ between adjacent bands' profiles, inverted). A boundary-slide analysis independently varied each boundary $\pm 25\%$ while holding others fixed.

\subsection{Within-Band Coordinate Tests}
\label{sec:within_methods}

Five tests evaluated within-band structure under the $(\fzero = 7.60, r = \phisym)$ system, using 200-bin Hz-corrected enrichment profiles.

\textbf{Test 1: Scaling comparison.} Five frequency scales (log, linear, ERB, mel, $\sqrt{\text{Hz}}$) were compared by polynomial $R^2$.

\textbf{Test 2: Landmark capture.} Voronoi $R^2$ of 12 $\phisym$-positions versus 12 equal-spaced, 12 rational, and 1000 random position sets.

\textbf{Test 3: Feature alignment.} Mean distance from enrichment extrema to nearest $\phisym$-position; 5000-permutation null.

\textbf{Test 4: Periodicity.} Autocorrelation of enrichment profiles at lags corresponding to $\phisym^{-k}$, $1/2$, $1/3$.

\textbf{Test 5: Noble vs.\ rational.} Enrichment at noble numbers versus simple rationals; 5000-permutation test.

\subsection{Voronoi Enrichment}
\label{sec:voronoi}

Twelve degree-6 positions per $\phisym$-octave were used as Voronoi landmarks:
\begin{equation}
\mathbf{u} = \{0,\; \phisym^{-6},\; \phisym^{-5},\; \phisym^{-4},\; \phisym^{-3},\; \phisym^{-2},\; 0.5,\; \phisym^{-1},\; 1-\phisym^{-3},\; 1-\phisym^{-4},\; 1-\phisym^{-5},\; 1-\phisym^{-6}\}.
\end{equation}

Each peak was assigned to its nearest position (circular distance). The boundary position ($u = 0$) was split into boundary\_lo and boundary\_hi for analysis of $\fzero$-convergence effects. Hz-weighted expected counts were computed per Equation~\ref{eq:hz_weighted_null}.

Six derived metrics summarise each band's profile: \textbf{mountain} (Noble$_1$ $-$ boundary), \textbf{ushape} (mean edges $-$ attractor), \textbf{peak\_height} (Noble$_1$ $-$ attractor), \textbf{ramp\_depth} (inv\_noble$_4$ $-$ noble$_4$), \textbf{center\_depletion} (attractor $-$ mean lower interior), and \textbf{asymmetry} (mean upper interior $-$ mean lower interior).

\subsection{Individual-Differences Analyses}
\label{sec:individual_diffs}

Per-subject enrichment profiles were computed for subjects with $\geq 30$ peaks per band after power filtering. Associations with continuous variables were quantified by Spearman $\rho$, corrected by Benjamini--Hochberg FDR ($q < 0.05$) within each analysis family. Age-partialed correlations used van der Waerden normal scores. Test--retest reliability (Dortmund longitudinal, $N = 208$) was estimated by ICC(2,1) (two-way random, absolute agreement, single measures).

\subsection{Computational Infrastructure}
\label{sec:infrastructure}

Peak extraction was parallelised on Google Cloud Platform (c2d-standard-32 VMs, custom image). Each dataset was processed as: VM spawn $\to$ GCS data copy $\to$ parallel per-subject extraction (28 cores) $\to$ results push $\to$ VM deletion ($\sim$5 minutes per dataset). The 48-configuration parameter sweep completed in $\sim$30 minutes. All code: Python 3.10 (NumPy 1.26, SciPy 1.11, pandas 2.1, specparam 1.0). All analysis code and extracted peak CSVs are archived on GCS.

\subsection{Ethics Statement}
\label{sec:ethics}

This study is a secondary analysis of publicly available, de-identified EEG datasets. All original studies obtained institutional ethics approval and informed consent from participants (or their guardians, in the case of HBN minors) prior to data collection and public release. No additional ethics approval was required for the present analyses of these de-identified data.

\subsection{Use of Artificial Intelligence Tools}
\label{sec:ai_disclosure}

Anthropic Claude Code (Opus 4.6) was used to assist with data analysis pipeline development and manuscript preparation. Anthropic Claude (Opus 4.6, chat) and OpenAI ChatGPT (5.4) were used as reviewers to check for consistency, clarity, errors, and omissions. All AI-generated content was reviewed, verified, and edited by the author, who takes full responsibility for the accuracy and integrity of the work.
